\section{Related Work}
\label{sec:related}

The task addressed in this paper lies at the intersection of fact-checking and information retrieval. A natural starting point is lexical retrieval, where probabilistic ranking models such as BM25 remain strong baselines thanks to their efficiency, interpretability and robustness in candidate generation \citep{RobertsonZaragoza2009}. In the setting of scientific claim source retrieval, lexical matching is especially useful when a post mentions recognizable entities or terminology that also appears in the title, abstract, venue or author list of the referenced paper.

Neural retrieval methods have been proposed to mitigate the lexical mismatch between queries and documents by learning dense semantic representations. More broadly, neural information retrieval has shifted attention from exact token overlap to representation learning and learned matching functions \citep{MarchesinEtAl2020}. In particular, dense dual-encoder architectures such as Dense Passage Retrieval (DPR) demonstrated that learned query and document embeddings can be highly effective for first-stage retrieval \citep{KarpukhinEtAl2020}. This line of work is especially relevant to our task because social media posts often paraphrase scientific findings instead of reusing the exact wording found in the original publication.

Between initial retrieval and final ranking, many modern systems adopt multi-stage pipelines that trade off efficiency and effectiveness. Late-interaction models such as ColBERT preserve fine-grained token-level matching while remaining substantially more scalable than fully pairwise neural scoring over the whole collection \citep{KhattabZaharia2020}. Likewise, transformer-based rerankers have become a common second-stage component for refining an initial candidate set with stronger query-document interactions \citep{NogueiraEtAl2020}. These approaches motivate retrieval architectures that first generate a manageable set of candidates and then apply more expensive relevance estimation only where it is most useful.



???????????????????????????
???????????????????????????
???????????????????????????
???????????????????????????
Se fine tuning va bene è da aggiungere parte su hard negatives e ricerca embedding incrociati usate:
https://arxiv.org/html/2505.18366v1
https://huggingface.co/blog/train-reranker

???????????????????????????
???????????????????????????
???????????????????????????
???????????????????????????











Our work is also related to the broader literature on automatic claim verification and, more specifically, to the CheckThat! shared tasks, which have addressed the detection, verification, and analysis of claims in social media across several editions \citep{ElsayedEtAl2019,BarronCedenoEtAl2020,NakovEtAl2021}; however, source retrieval for scientific web claims differs from standard verification tasks in an important way: the system is not asked to classify a claim as true or false, but to recover the original scientific paper implicitly referenced in the post and for this reason, the task is closely connected to evidence retrieval, while also requiring robust matching between informal multilingual posts and formal scientific documents.
